\documentclass[11pt]{amsart}

\usepackage[margin=1.15in]{geometry}
\usepackage{amsmath,amssymb,amsthm,mathtools}
\usepackage{enumitem}
\usepackage{microtype}
\usepackage[colorlinks=true,linkcolor=blue,citecolor=blue,urlcolor=blue]{hyperref}
\usepackage[nameinlink,capitalize]{cleveref}

\theoremstyle{plain}
\newtheorem{theorem}{Theorem}[section]
\newtheorem{proposition}[theorem]{Proposition}
\newtheorem{lemma}[theorem]{Lemma}
\newtheorem{corollary}[theorem]{Corollary}

\theoremstyle{definition}
\newtheorem{definition}[theorem]{Definition}

\theoremstyle{remark}
\newtheorem{remark}[theorem]{Remark}

\numberwithin{equation}{section}

\DeclareMathOperator{\diver}{div}

\title[Ancient Navier--Stokes criteria]
{Liouville and Lyapunov Criteria for Ancient Navier--Stokes Solutions}
\author{Guillem Duran Ballester}
\email{guillem@fragile.tech}
\date{}

\subjclass[2020]{35Q30, 35B40, 35B44, 35B45, 76D05}
\keywords{Navier--Stokes equations, ancient solutions, Type I blow-up, axisymmetric flows, Lyapunov functionals}

\begin{document}

\begin{abstract}
We study bounded ancient solutions of the three-dimensional Navier--Stokes
equations arising from Type I blow-up arguments.  We state standard PDE
hypotheses under which such solutions are spatially constant; in the Type I
class, these constants vanish.  We also prove the perturbative weighted
vorticity Lyapunov rigidity used by the unconditional Type I chain.

For the unforced equations, known axisymmetric Liouville theorems cover the
following cases: absence of swirl, pointwise scale-invariant control, finite
\(L_t^\infty L_x^p\) control of the swirl variable \(\Gamma=ru_\theta\), and
axial periodicity with bounded \(\Gamma\).  When one of these hypotheses holds
for a nonzero Type I blow-up limit, the Galilean-trivial conclusions of the
cited Liouville theorems reduce to the zero solution and contradict the
nontriviality of the limit.

We also record the Lyapunov rigidity argument coming from Gallay--Wayne's
small-data decay theorem for the three-dimensional rescaled vorticity equation.
It yields a coercive Lyapunov functional forcing any ancient trajectory that
stays in the perturbative weighted-vorticity ball to be zero.  The paper concludes
with a summary of the unconditional Liouville hypotheses and cited PDE results
used by the chain.
\end{abstract}

\maketitle

\section{Introduction}

\subsection{Ancient solutions from Type I blow-up}

The Type I blow-up argument in Paper I \cite{DuranPaperI} produces, after
parabolic rescaling around a Type I singular point, a nonzero ancient solution
\((U,P)\) of
\begin{equation}\label{eq:NS}
   \partial_tU+(U\cdot\nabla)U+\nabla P=\Delta U,
   \qquad \nabla\cdot U=0,
   \qquad t<0,
\end{equation}
satisfying the Type I ancient bound
\begin{equation}\label{eq:typeI}
   \sup_{t<0}\sqrt{-t}\,
   \|U(t)\|_{L^\infty(\mathbb R^3)}<\infty .
\end{equation}
Thus every result proving \(U\equiv0\) for an ancient solution satisfying
\eqref{eq:typeI} has the same consequence in the Type I problem:
\[
   \text{PDE hypothesis}
   \quad\Longrightarrow\quad
   U\equiv0
   \quad\Longrightarrow\quad
   \text{contradiction with nonzero blow-up limit}.
\]

The paper treats hypotheses for which a specific Liouville mechanism is
available.  The cases considered are:
\[
   \begin{gathered}
   \text{axisymmetric Liouville hypotheses},\qquad
   \text{small weighted rescaled vorticity}.
   \end{gathered}
\]
We make no assertion for decaying or symmetry-restricted ancient solutions
outside the hypotheses listed below.  Each conclusion is tied to an explicit
PDE hypothesis and a specified analytic result.

For the Type I application in Paper I, only those cases in which the stated
PDE hypothesis and the cited or proved Liouville theorem imply \(U\equiv0\)
are used as exclusions.  Ancient profiles outside these hypotheses are treated
separately.

\subsection{Axisymmetric Liouville hypotheses}

Several unforced axisymmetric cases are already covered by known Liouville
theorems.  Koch--Nadirashvili--Seregin--\v{S}ver\'ak cover the no-swirl case
and the pointwise scale-invariant condition \(|u|\le C/r\).  The
swirl-bearing finite-\(L^p\) condition
\[
   \Gamma=ru_\theta\in L_t^\infty L_x^p,\qquad 1\le p<\infty,
\]
is covered by Lei--Zhang--Zhao.  The \(z\)-periodic case with bounded
\(\Gamma\) is covered by Lei--Ren--Zhang.

The only additional step needed for Type I ancient limits is to convert
Galilean-triviality into zero in the blow-up frame.  We time-shift away from
the possible singular time \(t=0\), apply the bounded ancient Liouville
theorem, patch the resulting constants on backward intervals, and use
\eqref{eq:typeI}: if \(U(x,t)\equiv c\), then
\[
   \sup_{t<0}\sqrt{-t}\,\|U(t)\|_\infty
   =
   \sup_{t<0}\sqrt{-t}\,|c|
\]
is finite only when \(c=0\).

\subsection{Lyapunov hypotheses}

The Lyapunov result below uses a forward Lyapunov functional in the
Gallay--Wayne perturbative weighted-vorticity regime.  If \(w\)
solves the rescaled three-dimensional vorticity equation and remains in the
Gallay--Wayne small ball of \(L^2_\sigma(m)\), \(m>7/2\), then the functional
\begin{equation}\label{eq:intro-weighted-lyap}
   \mathcal L_m(w_0)
   =
   \sup_{\tau\ge0}e^{2\tau}\|\Phi_\tau w_0\|_m^2
\end{equation}
is coercive and contracts along the forward semiflow.  Sending the starting
time to \(-\infty\) then forces every ancient trajectory in the small ball to
be zero.

\subsection{Main results}

The first result records the unforced axisymmetric Liouville consequences
relevant to the Type I argument.

\begin{theorem}[Axisymmetric Liouville consequences for Type I ancient solutions]
\label{thm:axisymmetric-liouville}
Let \(U\) be a smooth Type I ancient solution of \eqref{eq:NS} satisfying
\eqref{eq:typeI}.  If \(U\) satisfies any of the following axisymmetric
assumptions, then \(U\equiv0\):
\begin{enumerate}[label=\textup{(\roman*)}]
\item axisymmetric without swirl;
\item axisymmetric and satisfying a pointwise scale-invariant bound
\[
   |U(x,t)|\le \frac{C}{r};
\]
\item axisymmetric with
\[
   \Gamma=rU_\theta\in L_t^\infty L_x^p
   \bigl(\mathbb R^3\times(-\infty,0)\bigr),
   \qquad 1\le p<\infty;
\]
\item axisymmetric, periodic in the physical axial variable \(z\), and with
\[
   \Gamma=rU_\theta\in L^\infty
   \bigl(\mathbb R^3\times(-\infty,0)\bigr).
\]
\end{enumerate}
Consequently none of these assumptions can hold for a nonzero Type I blow-up
limit.
\end{theorem}

\begin{theorem}[Perturbative weighted-vorticity Lyapunov rigidity]
\label{thm:weighted-vorticity-vanishing-intro}
Let \(m>\frac72\).  Let \(w(\tau)\) be an ancient mild solution of the
Gallay--Wayne rescaled vorticity equation
\begin{equation}\label{eq:weighted-vorticity-intro}
   \partial_\tau w
   =
   \Lambda w-(v\cdot\nabla)w+(w\cdot\nabla)v,
   \qquad \nabla\cdot w=0,
\end{equation}
with \(v\) recovered by the three-dimensional Biot--Savart law.  Assume
\[
   \sup_{\tau\in\mathbb R}\|w(\tau)\|_{L^2_\sigma(m)}
   \le \rho_m,
\]
where \(\rho_m\) is the Gallay--Wayne small-data radius.  Then
\[
   w\equiv0.
\]
\end{theorem}

Here ``ancient mild'' means that every forward time-shift of \(w\) is the
Gallay--Wayne mild solution launched from its value at the starting time.

\begin{theorem}[Consequences for Type I blow-up limits]
\label{thm:typeI-consequences}
Let \(U\) be a nonzero Type I blow-up limit.  Then \(U\) cannot satisfy any of
the axisymmetric assumptions in \Cref{thm:axisymmetric-liouville}; nor can its
corresponding rescaled vorticity satisfy the hypotheses of
\Cref{thm:weighted-vorticity-vanishing-intro}.
\end{theorem}

Only the explicit hypotheses stated above are used.

\subsection{Relation to Papers I and II}

Paper I \cite{DuranPaperI} reduces Type I singularities to nonzero ancient
limits satisfying \eqref{eq:typeI}.  Paper II \cite{DuranPaperII} treats
uniformly \(L^3\)-tight ancient solutions.  The present paper records the
unconditional symmetry and perturbative weighted-vorticity exclusions used by
the chain.  Known Liouville theorems rule out the axisymmetric cases listed
above, and the Lyapunov argument handles the perturbative weighted-vorticity
case.
Paper III does not depend on Paper II.

\subsection{Summary of hypotheses}

\begin{center}
\begin{tabular}{p{0.36\textwidth}p{0.18\textwidth}p{0.34\textwidth}}
\hline
Hypothesis & Conclusion & Analytic ingredient \\
\hline
axisymmetric no swirl & \(U\equiv0\) & KNSS \\
pointwise \(|u|\le C/r\) & \(U\equiv0\) & KNSS \\
\(\Gamma\in L_t^\infty L_x^p\), \(1\le p<\infty\) & \(U\equiv0\) & LZZ \\
\(z\)-periodic bounded \(\Gamma\) & \(U\equiv0\) & LRZ \\
small weighted vorticity & \(w\equiv0\) & Gallay--Wayne \\
\hline
\end{tabular}
\end{center}

\subsection{Organization}

\Cref{sec:typeI} records the Type I ancient setup, self-similar variables, the
invariance facts, and the time-shift lemma.
\Cref{sec:axisymmetric-liouville-theorems} fixes axisymmetric notation and the
known Liouville theorems.  \Cref{sec:axisymmetric-consequences} proves the
unforced axisymmetric Liouville consequences.
\Cref{sec:weighted-vorticity} proves the Gallay--Wayne Lyapunov rigidity, and
\Cref{sec:summary} summarizes the consequences.

\section{Ancient Type I limits and self-similar variables}
\label{sec:typeI}

\subsection{Type I ancient solutions}

\begin{definition}[Type I ancient solution]
\label{def:typeI-ancient}
A smooth mild solution \(U\) of \eqref{eq:NS} on
\(\mathbb R^3\times(-\infty,0)\) is called Type I ancient if
\eqref{eq:typeI} holds.  Mild is understood in the standard velocity
formulation \cite{Kato1984,Sohr2001}: for every \(-\infty<s<t<0\),
\begin{equation}\label{eq:mild-formula}
   U(t)=e^{(t-s)\Delta}U(s)
   -\int_s^t e^{(t-\sigma)\Delta}\mathbb P\nabla\cdot
      (U(\sigma)\otimes U(\sigma))\,d\sigma
\end{equation}
in distributions, where \(\mathbb P\) is the Leray projector.
\end{definition}

\subsection{Nonzero Type I blow-up limits}

\begin{definition}[Nonzero Type I blow-up limit]
\label{def:nonzero-typeI-limit}
A Type I ancient solution \(U\) is called a nonzero Type I blow-up limit if it
arises as the limit of parabolic rescalings around a Type I singular point and
the limiting solution is not identically zero.  The arguments below require
only the conclusion \(U\not\equiv0\).
\end{definition}

\begin{lemma}[Nontriviality of blow-up limits]
\label{lem:nonzero-typeI-limit}
   A nonzero Type I blow-up limit is not identically zero.
\end{lemma}

\begin{proof}
This is part of the definition above.
\end{proof}

\subsection{Backward self-similar variables}

For \(t<0\), define
\begin{equation}\label{eq:self-similar}
   V(y,\tau)=\sqrt{-t}\,U(x,t),\qquad
   Q(y,\tau)=(-t)P(x,t),\qquad
   y=\frac{x}{\sqrt{-t}},\qquad
   \tau=-\log(-t).
\end{equation}
Equivalently,
\[
   x=e^{-\tau/2}y,\qquad t=-e^{-\tau}.
\]

\begin{lemma}[Self-similar equation]\label{lem:self-similar-equation}
Let \((U,P)\) be a smooth solution of \eqref{eq:NS}, and define \((V,Q)\) by
\eqref{eq:self-similar}.  Then \(V\) is defined on
\(\mathbb R^3\times\mathbb R\) and satisfies
\begin{equation}\label{eq:centered}
   \partial_\tau V+(V\cdot\nabla)V+\nabla Q
   =
   \Delta V-\frac12(V+y\cdot\nabla V),
   \qquad \nabla\cdot V=0.
\end{equation}
Moreover \eqref{eq:typeI} is equivalent to
\[
   \sup_{\tau\in\mathbb R}
   \|V(\tau)\|_{L^\infty(\mathbb R^3)}<\infty .
\]
\end{lemma}

\begin{proof}
Let \(a=\sqrt{-t}=e^{-\tau/2}\).  Then \(V=aU\), \(y=x/a\), and
\[
   \partial_t\tau=\frac1{-t}=\frac1{a^2},
   \qquad
   \partial_t y=\frac{y}{2a^2},
   \qquad
   \partial_t(a^{-1})=\frac1{2a^3}.
\]
Since \(U=a^{-1}V\), differentiating at fixed \(x\) gives
\[
   \partial_tU
   =
   a^{-3}
   \left(
      \partial_\tau V+\frac12 y\cdot\nabla V+\frac12V
   \right).
\]
Spatial derivatives scale as
\[
   \nabla_xU=a^{-2}\nabla_yV,\qquad
   \Delta_xU=a^{-3}\Delta_yV,
\]
and the pressure scaling gives
\[
   \nabla_xP=a^{-3}\nabla_yQ.
\]
Thus
\[
   (U\cdot\nabla_x)U=a^{-3}(V\cdot\nabla_y)V.
\]
Multiplying \eqref{eq:NS} by \(a^3\) yields \eqref{eq:centered}; the
divergence equation follows from
\(\nabla_x\cdot U=a^{-2}\nabla_y\cdot V\).  The spatial map
\(y\mapsto x=e^{-\tau/2}y\) is a bijection for each \(\tau\), and
\[
   \|V(\tau)\|_\infty
   =
   \sqrt{-t}\,\|U(t)\|_\infty .
\]
Taking suprema over \(t<0\), equivalently \(\tau\in\mathbb R\), gives the
last assertion.
\end{proof}

\subsection{Preservation of axisymmetry and swirl variables}

For axisymmetric flows we write
\[
   x=(r\cos\theta,r\sin\theta,z),\qquad
   U=U_re_r+U_\theta e_\theta+U_ze_z,
\]
and in centered variables
\[
   y=(\rho\cos\theta,\rho\sin\theta,Z),\qquad
   V=V_\rho e_\rho+V_\theta e_\theta+V_Ze_Z.
\]
The scalar swirl variables are
\[
   \Gamma_U(r,z,t)=rU_\theta(r,z,t),\qquad
   \Gamma_V(\rho,Z,\tau)=\rho V_\theta(\rho,Z,\tau).
\]

\begin{lemma}[Preservation of axisymmetric structure]
\label{lem:structure-preserved}
Axisymmetry and absence of swirl are preserved under
\eqref{eq:self-similar}.  If \(U\) is axisymmetric, then
\[
   \Gamma_V(\rho,Z,\tau)=\Gamma_U(r,z,t),
   \qquad
   (r,z)=\sqrt{-t}\,(\rho,Z).
\]
Fixed physical \(z\)-periodicity is preserved by physical time shifts, but
under the centered change of variables it becomes a \(\tau\)-dependent period
in the \(Z\)-coordinate.
\end{lemma}

\begin{proof}
The self-similar map is a scalar spatial dilation and a scalar multiplication
of the velocity, so it commutes with rotations around the symmetry axis.
The cylindrical basis is unchanged by radial dilation, and
\[
   V_\theta(\rho,Z,\tau)
   =
   \sqrt{-t}\,U_\theta(r,z,t),
   \qquad
   r=\sqrt{-t}\rho,\quad z=\sqrt{-t}Z .
\]
Thus \(U_\theta\equiv0\) if and only if \(V_\theta\equiv0\), and
\[
   \Gamma_V=\rho V_\theta
   =
   \rho\sqrt{-t}\,U_\theta
   =
   rU_\theta
   =
   \Gamma_U .
\]
If \(U\) has physical axial period \(L\), then \(V\) has period
\(L/\sqrt{-t}=Le^{\tau/2}\) in \(Z\), not a fixed period.
\end{proof}

\subsection{Time shifts away from the singular time}

\begin{lemma}[Time shifts away from the singular time]
\label{lem:time-shift}
Let \(U\) be a smooth Type I ancient solution and fix \(T<0\).  Then
\[
   W_T(x,s)=U(x,T+s),\qquad \Pi_T(x,s)=P(x,T+s),\qquad s<0,
\]
is a smooth bounded mild ancient solution of \eqref{eq:NS} on
\(\mathbb R^3\times(-\infty,0)\).  Moreover the axisymmetric, no-swirl,
\(L^p\)-swirl, bounded-\(\Gamma\), and physical \(z\)-periodic properties are
preserved.
\end{lemma}

\begin{proof}
Time translation preserves the differential equation because
\(\partial_sW_T(x,s)=\partial_tU(x,T+s)\) and all spatial derivatives are
unchanged.  For the mild formulation, apply \eqref{eq:mild-formula} to \(U\)
between the physical times \(T+s<T+t<0\):
\[
   U(T+t)=e^{(t-s)\Delta}U(T+s)
   -\int_{T+s}^{T+t}e^{(T+t-\eta)\Delta}\mathbb P\nabla\cdot
      (U(\eta)\otimes U(\eta))\,d\eta .
\]
With \(\eta=T+\sigma\), this becomes the mild formula for \(W_T\) on the
interval \((s,t)\).  For \(s<0\), \(T+s<T<0\), so the Type I bound gives
\[
   \|W_T(s)\|_\infty
   =
   \|U(T+s)\|_\infty
   \le
   \frac1{\sqrt{-T}}
   \sup_{\eta<0}\sqrt{-\eta}\,\|U(\eta)\|_\infty .
\]
The invariance assertions follow because only the time variable has been
translated.  In particular
\[
   \Gamma_{W_T}(r,z,s)=\Gamma_U(r,z,T+s),
\]
so \(L_s^\infty L_x^p\) and \(L^\infty\) bounds for \(\Gamma\) restrict to the
shifted solution, and a fixed physical axial period remains the same.
\end{proof}

\subsection{PDE hypotheses used below}

The arguments below are stated directly in terms of hypotheses on ancient
solutions.  The unconditional axisymmetric hypotheses are:
\begin{enumerate}[label=\textup{(\roman*)}]
\item axisymmetry and \(U_\theta\equiv0\);
\item axisymmetry and the pointwise bound \(|U(x,t)|\le C/r\) for \(r>0\);
\item axisymmetry and
\[
   \Gamma=rU_\theta\in L_t^\infty L_x^p
   \bigl(\mathbb R^3\times(-\infty,0)\bigr),
   \qquad 1\le p<\infty;
\]
\item axisymmetry, periodicity in the physical axial variable \(z\), and
\(\Gamma\in L^\infty\).
\end{enumerate}
The additional hypothesis used later is the Gallay--Wayne small
weighted-vorticity condition.

\section{Axisymmetric notation and known Liouville theorems}
\label{sec:axisymmetric-liouville-theorems}

\subsection{Cylindrical variables and the swirl scalar}

For an axisymmetric solution \(u\), write
\[
   u=u_re_r+u_\theta e_\theta+u_ze_z,\qquad
   \Gamma=ru_\theta,\qquad
   b=u_re_r+u_ze_z .
\]
When \(\Gamma\) is placed in \(L^p(\mathbb R^3)\), it is regarded as the
corresponding function of \(x=(r\cos\theta,r\sin\theta,z)\), independent of
\(\theta\):
\[
   \|\Gamma(t)\|_{L^p(\mathbb R^3)}^p
   =
   2\pi\int_{\mathbb R}\int_0^\infty
      |\Gamma(r,z,t)|^p r\,dr\,dz,\qquad 1\le p<\infty .
\]

\subsection{Constant and Galilean-trivial solutions}

\begin{definition}[Galilean-trivial]
A velocity field is Galilean-trivial if, after a Galilean transform, it is a
constant vector field.  In the Type I blow-up frame below we use only the
following consequence: when a cited Liouville theorem gives a constant
velocity field in that frame, the Type I bound forces the constant to be zero.
\end{definition}

\begin{lemma}[Spatially constant mild solutions]\label{lem:spatial-constant-mild}
Let \(W\) be a smooth mild solution of \eqref{eq:NS} on
\(\mathbb R^3\times(-\infty,0)\).  If \(W(x,t)=b(t)\) for all \(x\) and \(t\),
then \(b\) is constant in time.
\end{lemma}

\begin{proof}
Insert \(W(x,t)=b(t)\) into \eqref{eq:mild-formula}.  The heat semigroup
leaves constant vector fields fixed, and
\(\nabla\cdot(b(\sigma)\otimes b(\sigma))=0\).  Thus
\(b(t)=b(s)\) for every \(s<t<0\).
\end{proof}

\begin{lemma}[Patching constants on backward intervals]
\label{lem:patch-constants}
Let \(U:\mathbb R^3\times(-\infty,0)\to\mathbb R^3\) be a function.  Assume
that for every \(T<0\) there exists \(c_T\in\mathbb R^3\) such that
\[
   U(x,t)=c_T\qquad(x\in\mathbb R^3,\ t<T).
\]
Then there is a single \(c\in\mathbb R^3\) such that
\[
   U(x,t)=c\qquad(x\in\mathbb R^3,\ t<0).
\]
\end{lemma}

\begin{proof}
If \(T_1<T_2<0\), the two constants agree on the nonempty interval
\((-\infty,T_1)\).  Hence all \(c_T\) are equal.  For any \(t<0\), choose
\(T\) with \(t<T<0\).
\end{proof}

\begin{lemma}[Constant Type I ancient solutions vanish]
\label{lem:galilean-typeI-zero}
Let \(U\) be a Type I ancient solution.  If \(U(x,t)\equiv c\), then
\(c=0\).
\end{lemma}

\begin{proof}
The Type I bound gives
\[
   \sqrt{-t}\,|c|
   \le
   \sup_{s<0}\sqrt{-s}\,\|U(s)\|_\infty<\infty
   \qquad(t<0).
\]
Letting \(t\to-\infty\) gives \(c=0\).
\end{proof}

\subsection{Known Liouville theorems}

We use the following Liouville theorems in the stated forms.

\begin{theorem}[KNSS no-swirl Liouville theorem]
\label{thm:knss-noswirl}
Every bounded ancient mild axisymmetric no-swirl solution of the unforced
Navier--Stokes equations is Galilean-trivial.  In the form used below, the
solution is a spatially constant axial drift.
\end{theorem}

\begin{proof}
This is the no-swirl Liouville theorem of
Koch--Nadirashvili--Seregin--\v{S}ver\'ak \cite{KNSS2009}, stated in the form
needed for the Type I vanishing argument.
\end{proof}

\begin{theorem}[KNSS pointwise scale-invariant condition]
\label{thm:knss-pointwise}
Every bounded ancient mild axisymmetric solution satisfying
\[
   |u(x,t)|\le \frac{C}{r}
   \qquad(r>0,\ t<0)
\]
is Galilean-trivial.  In the form used below, it is identically zero.
\end{theorem}

\begin{proof}
This is the pointwise Liouville theorem of \cite{KNSS2009}, stated in the form
used here.
\end{proof}

\begin{theorem}[Lei--Zhang--Zhao]
\label{thm:lzz}
Let \(u\) be a bounded ancient mild axisymmetric solution.  If
\[
   \Gamma=ru_\theta\in L_t^\infty L_x^p
   \bigl(\mathbb R^3\times(-\infty,0)\bigr),
   \qquad 1\le p<\infty,
\]
then \(u\) is a constant vector field.
\end{theorem}

\begin{proof}
This is the ancient axisymmetric \(L^p\)-swirl Liouville theorem of
Lei--Zhang--Zhao \cite{LZZ}.
\end{proof}

\begin{theorem}[Lei--Ren--Zhang]
\label{thm:lrz}
Let \(u\) be a bounded ancient mild axisymmetric solution, periodic in \(z\),
with
\[
   \Gamma=ru_\theta\in L^\infty
   \bigl(\mathbb R^3\times(-\infty,0)\bigr).
\]
Then \(u\) is a constant axial vector field.
\end{theorem}

\begin{proof}
This is the periodic bounded-swirl theorem of Lei--Ren--Zhang
\cite{LRZ}.  Their statement is on \(\mathbb R^2\times\mathbb T^1\); a
physical \(z\)-periodic whole-space solution descends to the quotient, and a
Navier--Stokes scaling normalizes the period.
\end{proof}

\section{Unforced axisymmetric Liouville consequences}
\label{sec:axisymmetric-consequences}

\subsection{No swirl}

\begin{proposition}[No-swirl Type I vanishing]\label{prop:noswirl-typeI}
Let \(U\) be a smooth Type I ancient solution.  Assume \(U\) is axisymmetric
and \(U_\theta\equiv0\).  Then \(U\equiv0\).
\end{proposition}

\begin{proof}
Fix \(T<0\) and set \(W_T(x,s)=U(x,T+s)\).  By \Cref{lem:time-shift},
\(W_T\) is a bounded ancient mild axisymmetric no-swirl solution.  Applying
\Cref{thm:knss-noswirl} gives that \(W_T\) is a spatially constant axial
drift, say \(W_T(x,s)=b_T(s)e_z\).  By
\Cref{lem:spatial-constant-mild}, \(b_T(s)\) is constant in \(s\).  Thus
\(U(x,t)=c_T\) for \(t<T\).  Since \(T<0\) is arbitrary,
\Cref{lem:patch-constants} gives \(U(x,t)=c\) for all \(t<0\).  Finally,
\Cref{lem:galilean-typeI-zero} gives \(c=0\).
\end{proof}

\subsection{Pointwise scale-invariant control}

\begin{proposition}[Pointwise \(1/r\) Type I vanishing]\label{prop:pointwise-typeI}
Let \(U\) be a smooth Type I ancient solution.  Assume \(U\) is axisymmetric
and satisfies
\[
   |U(x,t)|\le \frac{C}{r}
   \qquad(r>0,\ t<0).
\]
Then \(U\equiv0\).
\end{proposition}

\begin{proof}
Fix \(T<0\).  The shifted solution \(W_T(x,s)=U(x,T+s)\) is bounded, ancient,
axisymmetric, and satisfies the same pointwise bound \(|W_T|\le C/r\).
\Cref{thm:knss-pointwise} gives \(W_T\equiv0\).  Hence \(U\equiv0\) on
\((-\infty,T)\).  Since \(T<0\) is arbitrary, \(U\equiv0\).
\end{proof}

\subsection{\texorpdfstring{Finite \(L^p\) control of \(\Gamma\)}
{Finite Lp control of Gamma}}

\begin{proposition}[Finite-\(L^p\) swirl Type I vanishing]
\label{prop:lp-swirl-typeI}
Let \(U\) be a smooth Type I ancient solution.  Assume \(U\) is axisymmetric
and, for some \(1\le p<\infty\),
\[
   \Gamma=rU_\theta\in L_t^\infty L_x^p
   \bigl(\mathbb R^3\times(-\infty,0)\bigr).
\]
Then \(U\equiv0\).
\end{proposition}

\begin{proof}
Fix \(T<0\).  By \Cref{lem:time-shift}, \(W_T(x,s)=U(x,T+s)\) is a bounded
ancient mild axisymmetric solution and
\[
   \Gamma_{W_T}\in L_s^\infty L_x^p
   \bigl(\mathbb R^3\times(-\infty,0)\bigr).
\]
\Cref{thm:lzz} gives that \(W_T\) is a constant vector field.  Patch these
constants in \(T\) by \Cref{lem:patch-constants} and then use
\Cref{lem:galilean-typeI-zero}.
\end{proof}

\subsection{Periodic bounded swirl}

\begin{proposition}[Periodic bounded-swirl Type I vanishing]
\label{prop:periodic-swirl-typeI}
Let \(U\) be a smooth Type I ancient solution.  Assume \(U\) is axisymmetric,
periodic in the physical axial variable \(z\), and satisfies
\[
   \Gamma=rU_\theta\in L^\infty
   \bigl(\mathbb R^3\times(-\infty,0)\bigr).
\]
Then \(U\equiv0\).
\end{proposition}

\begin{proof}
Fix \(T<0\).  By \Cref{lem:time-shift}, the shifted solution \(W_T\) is
bounded, ancient, axisymmetric, physically \(z\)-periodic, and has bounded
\(\Gamma\).  \Cref{thm:lrz} gives that \(W_T\) is a constant axial vector
field.  The constants patch as \(T\) varies, and the Type I bound eliminates
the resulting constant by \Cref{lem:galilean-typeI-zero}.
\end{proof}

\subsection{Proof of the axisymmetric Liouville theorem}

\begin{proof}[Proof of \Cref{thm:axisymmetric-liouville}]
The four cases are precisely the hypotheses of
\Cref{prop:noswirl-typeI,prop:pointwise-typeI,prop:lp-swirl-typeI,prop:periodic-swirl-typeI}.
Each proposition gives \(U\equiv0\).  This contradicts
\Cref{lem:nonzero-typeI-limit} when \(U\) is a nonzero Type I blow-up limit.
\end{proof}

\section{Perturbative weighted-vorticity Lyapunov rigidity}
\label{sec:weighted-vorticity}

\subsection{The rescaled vorticity equation}

In this section \(\tau\) is the forward self-similar time used in the
Gallay--Wayne vorticity formulation.  The ancient trajectories are indexed by
all \(\tau\in\mathbb R\), and the forward semiflow acts by increasing
\(\tau\).

Let \(w=\nabla\times v\).  The rescaled vorticity equation is
\begin{equation}\label{eq:weighted-vorticity}
   \partial_\tau w
   =
   \Lambda w-(v\cdot\nabla)w+(w\cdot\nabla)v,
   \qquad \nabla\cdot w=0,
\end{equation}
where
\begin{equation}\label{eq:lambda}
   \Lambda w=\Delta w+\frac12\xi\cdot\nabla w+w,
\end{equation}
and \(v\) is recovered from \(w\) by the three-dimensional Biot--Savart law.
For \(m\ge0\), set
\[
   L^2_\sigma(m)
   =
   \left\{
      f\in L^2(\mathbb R^3;\mathbb R^3):
      \nabla\cdot f=0,\quad \|f\|_m<\infty
   \right\},
\]
where
\[
   \|f\|_m^2
   =
   \int_{\mathbb R^3}(1+|\xi|)^{2m}|f(\xi)|^2\,d\xi .
\]

\subsection{Gallay--Wayne small-data decay}

\begin{theorem}[Gallay--Wayne small-data decay]\label{thm:gallay}
Let \(0<\mu\le1\) and \(m>2\mu+\frac12\).  There exist constants
\(\rho_m>0\) and \(C_m\ge1\) such that if \(\|w_0\|_m\le\rho_m\), then
\eqref{eq:weighted-vorticity} has a unique global mild solution
\(w(\tau)=\Phi_\tau w_0\) in \(C([0,\infty);L^2_\sigma(m))\) and
\[
   \|\Phi_\tau w_0\|_m
   \le
   C_m e^{-\mu\tau}\|w_0\|_m,\qquad \tau\ge0.
\]
In particular, if \(m>\frac72\), one may use \(\mu=1\).
\end{theorem}

\begin{proof}
This is the Gallay--Wayne weighted small-data theorem for the rescaled
three-dimensional vorticity equation \cite{GallayWayne2002}.  The only
conclusions used below are existence, uniqueness, the semiflow property, and
the displayed exponential decay.
\end{proof}

\subsection{The Lyapunov functional}

\begin{definition}[Perturbative weighted-vorticity condition]
\label{def:weighted-vorticity-condition}
Fix \(m>\frac72\), and let \(\rho_m\) be the radius in \Cref{thm:gallay}.
An ancient vorticity trajectory satisfies the perturbative
weighted-vorticity condition with exponent \(m\) if it is a function
\[
   w\in C(\mathbb R;L^2_\sigma(m))
\]
that solves \eqref{eq:weighted-vorticity} for all \(\tau\in\mathbb R\) and
satisfies
\[
   \sup_{\tau\in\mathbb R}\|w(\tau)\|_m\le\rho_m .
\]
Solving for all \(\tau\) means that every time shift is compatible with the
Gallay--Wayne forward mild semiflow.
\end{definition}

\begin{proposition}[Semiflow Lyapunov functional]\label{prop:weighted-lyap}
Fix \(m>\frac72\).  For \(\|w_0\|_m\le\rho_m\), define
\begin{equation}\label{eq:weighted-lyap}
   \mathcal L_m(w_0)
   =
   \sup_{\tau\ge0}e^{2\tau}\|\Phi_\tau w_0\|_m^2 .
\end{equation}
Then \(\mathcal L_m\) is well-defined and satisfies
\begin{equation}\label{eq:weighted-coercive}
   \|w_0\|_m^2
   \le
   \mathcal L_m(w_0)
   \le
   C_m^2\|w_0\|_m^2 .
\end{equation}
Moreover, for every \(s\ge0\),
\begin{equation}\label{eq:weighted-contract}
   \mathcal L_m(\Phi_s w_0)
   \le
   e^{-2s}\mathcal L_m(w_0).
\end{equation}
\end{proposition}

\begin{proof}
The Gallay--Wayne estimate with \(\mu=1\) gives
\[
   e^{2\tau}\|\Phi_\tau w_0\|_m^2
   \le C_m^2\|w_0\|_m^2,
\]
which proves finiteness and the upper bound.  The lower bound is the value at
\(\tau=0\).  By uniqueness,
\[
   \Phi_\tau(\Phi_s w_0)=\Phi_{\tau+s}w_0 .
\]
Hence
\[
   \mathcal L_m(\Phi_s w_0)
   =
   \sup_{\tau\ge0}e^{2\tau}\|\Phi_{\tau+s}w_0\|_m^2
   =
   e^{-2s}\sup_{\sigma\ge s}e^{2\sigma}\|\Phi_\sigma w_0\|_m^2
   \le e^{-2s}\mathcal L_m(w_0).
\]
\end{proof}

\subsection{Ancient rigidity in the small weighted ball}

\begin{proof}[Proof of \Cref{thm:weighted-vorticity-vanishing-intro}]
Let \(w\) be an ancient solution in the small ball and fix
\(\tau_0\in\mathbb R\).  For \(s<\tau_0\), the shifted function
\[
   W_s(\theta)=w(s+\theta),\qquad \theta\ge0,
\]
is the Gallay--Wayne mild solution with initial datum \(w(s)\).  The small-ball
hypothesis gives \(\|w(s)\|_m\le\rho_m\), so the forward solution is in the
domain of \(\Phi_\theta\), and uniqueness gives
\[
   W_s(\theta)=\Phi_\theta w(s),\qquad \theta\ge0.
\]
Choosing \(\theta=\tau_0-s\) gives
\[
   w(\tau_0)=\Phi_{\tau_0-s}w(s).
\]
Since \(\|w(\tau_0)\|_m\le\rho_m\), the Lyapunov functional is also defined at
\(w(\tau_0)\).  Applying \Cref{prop:weighted-lyap} with forward time
\(\tau_0-s\),
\[
   \mathcal L_m(w(\tau_0))
   \le
   e^{-2(\tau_0-s)}\mathcal L_m(w(s))
   \le
   C_m^2\rho_m^2e^{-2(\tau_0-s)} .
\]
Letting \(s\to-\infty\) gives \(\mathcal L_m(w(\tau_0))=0\).  The coercive
lower bound \eqref{eq:weighted-coercive} gives \(w(\tau_0)=0\).  Since
\(\tau_0\) was arbitrary, \(w\equiv0\).
\end{proof}

\section{Summary of PDE consequences}
\label{sec:summary}

\subsection{Zero conclusions}

\begin{theorem}[Zero conclusions for Type I blow-up limits]
\label{thm:typeI-zero-hypotheses}
The following hypotheses cannot hold for a nonzero Type I blow-up limit:
\begin{enumerate}[label=\textup{(\roman*)}]
\item axisymmetry and \(U_\theta\equiv0\);
\item axisymmetry and the pointwise bound \(|U(x,t)|\le C/r\);
\item axisymmetry and finite-\(L^p\) control of the swirl scalar,
\[
   \Gamma\in L_t^\infty L_x^p,\qquad 1\le p<\infty;
\]
\item axisymmetry, \(z\)-periodicity, and bounded \(\Gamma\);
\item the Gallay--Wayne small weighted-vorticity condition, whenever the
velocity is recovered from the corresponding vorticity by the Biot--Savart law.
\end{enumerate}
\end{theorem}

\begin{proof}
Items (i)--(iv) are \Cref{thm:axisymmetric-liouville}.  Item (v) is
\Cref{thm:weighted-vorticity-vanishing-intro}; since the velocity in that
case is recovered from \(w\) by the Biot--Savart law, \(w\equiv0\) gives the
zero velocity.  In each Type I case, the zero conclusion contradicts
\Cref{lem:nonzero-typeI-limit}.
\end{proof}

\subsection{Consequences for Type I blow-up arguments}

\begin{proof}[Proof of \Cref{thm:typeI-consequences}]
The axisymmetric cases are ruled out by \Cref{thm:axisymmetric-liouville}.  The
perturbative weighted-vorticity case is ruled out by
\Cref{thm:weighted-vorticity-vanishing-intro} once the corresponding rescaled
vorticity satisfies the small weighted condition; the Biot--Savart
reconstruction then gives zero velocity.  These zero conclusions contradict
the nontriviality of a Type I blow-up limit.
\end{proof}

\begin{theorem}[Combined zero conclusion]
\label{thm:stated-hypotheses-zero}
Let \(U\) be a Type I ancient solution obtained as a nonzero Type I blow-up
limit.  If \(U\) satisfies one of the hypotheses in
\Cref{thm:typeI-zero-hypotheses}, then \(U\equiv0\), a contradiction.
\end{theorem}

\begin{proof}
This is \Cref{thm:typeI-zero-hypotheses}.
\end{proof}

\subsection{Summary table}

\begin{center}
\begin{tabular}{p{0.36\textwidth}p{0.18\textwidth}p{0.34\textwidth}}
\hline
Hypothesis & Conclusion & Analytic ingredient \\
\hline
no swirl & \(U\equiv0\) & KNSS plus Type I constant removal \\
pointwise \(1/r\) & \(U\equiv0\) & KNSS \\
\(\Gamma\in L_t^\infty L_x^p\) & \(U\equiv0\) & LZZ plus Type I constant removal \\
periodic bounded \(\Gamma\) & \(U\equiv0\) & LRZ plus Type I constant removal \\
small weighted vorticity & \(w\equiv0\) & Gallay--Wayne Lyapunov functional \\
\hline
\end{tabular}
\end{center}

The no-swirl, pointwise scale-invariant, finite-\(L^p\)-swirl, and periodic
bounded-swirl cases are covered by known axisymmetric Liouville theorems.  The
perturbative weighted-vorticity case is covered by a Lyapunov functional
derived from Gallay--Wayne decay.
Thus every zero conclusion used in the Type I application rests on an explicit
PDE hypothesis stated in this paper.

\begin{thebibliography}{99}

\bibitem{CSTY2008}
C.-C. Chen, R. M. Strain, T.-P. Tsai, and H.-T. Yau,
\emph{Lower bound on the blow-up rate of the axisymmetric Navier--Stokes
equations},
Int. Math. Res. Not. IMRN \textbf{2008}, Art. ID rnn016, 31 pp.,
doi:10.1093/imrn/rnn016.

\bibitem{CSTY2009}
C.-C. Chen, R. M. Strain, T.-P. Tsai, and H.-T. Yau,
\emph{Lower bounds on the blow-up rate of the axisymmetric Navier--Stokes
equations II},
Comm. Partial Differential Equations \textbf{34} (2009), no. 3, 203--232,
doi:10.1080/03605300902793956.

\bibitem{DuranPaperI}
G. Duran Ballester,
\emph{Type I blow-up limits and a remaining-case criterion for the
three-dimensional Navier--Stokes equations}, preprint, 2026.

\bibitem{DuranPaperII}
G. Duran Ballester,
\emph{Uniformly tight ancient Navier--Stokes dynamics and endpoint \(L^3\)
rigidity}, preprint, 2026.

\bibitem{ESS2003}
L. Escauriaza, G. A. Seregin, and V. \v{S}ver\'ak,
\emph{\(L_{3,\infty}\)-solutions of the Navier--Stokes equations and backward
uniqueness},
Russian Math. Surveys \textbf{58} (2003), no. 2, 211--250,
doi:10.1070/RM2003v058n02ABEH000609.

\bibitem{GallayWayne2002}
T. Gallay and C. E. Wayne,
\emph{Long-time asymptotics of the Navier--Stokes and vorticity equations on
\(\mathbb R^3\)},
Philos. Trans. Roy. Soc. London Ser. A \textbf{360} (2002), no. 1799,
2155--2188, doi:10.1098/rsta.2002.1068.

\bibitem{Kato1984}
T. Kato,
\emph{Strong \(L^p\)-solutions of the Navier--Stokes equation in
\(\mathbb R^m\), with applications to weak solutions},
Math. Z. \textbf{187} (1984), no. 4, 471--480,
doi:10.1007/BF01174182.

\bibitem{KNSS2009}
G. Koch, N. Nadirashvili, G. Seregin, and V. \v{S}ver\'ak,
\emph{Liouville theorems for the Navier--Stokes equations and applications},
Acta Math. \textbf{203} (2009), no. 1, 83--105,
doi:10.1007/s11511-009-0039-6.

\bibitem{LeiZhang2011}
Z. Lei and Q. S. Zhang,
\emph{A Liouville theorem for the axially-symmetric Navier--Stokes equations},
J. Funct. Anal. \textbf{261} (2011), no. 8, 2323--2345,
doi:10.1016/j.jfa.2011.06.016.

\bibitem{LRZ}
Z. Lei, X. Ren, and Q. S. Zhang,
\emph{A Liouville theorem for axially symmetric Navier--Stokes equations on
\(\mathbb R^2\times\mathbb T^1\)},
Math. Ann. \textbf{383} (2022), no. 1--2, 415--431,
doi:10.1007/s00208-020-02128-9.

\bibitem{LZZ}
Z. Lei, Q. Zhang, and N. Zhao,
\emph{Improved Liouville theorems for axially symmetric Navier--Stokes
equations},
Sci. Sin. Math. \textbf{47} (2017), no. 10, 1183--1198,
doi:10.1360/N012016-00149.

\bibitem{Seregin2012}
G. Seregin,
\emph{A certain necessary condition of potential blow up for Navier--Stokes
equations},
Comm. Math. Phys. \textbf{312} (2012), no. 3, 833--845,
doi:10.1007/s00220-011-1391-x.

\bibitem{SereginSverak2009}
G. Seregin and V. \v{S}ver\'ak,
\emph{On Type I singularities of the local axi-symmetric solutions of the
Navier--Stokes equations},
Comm. Partial Differential Equations \textbf{34} (2009), no. 2, 171--201,
doi:10.1080/03605300802683687.

\bibitem{Sohr2001}
H. Sohr,
\emph{The Navier--Stokes Equations: An Elementary Functional Analytic
Approach},
Birkh\"auser, Basel, 2001,
doi:10.1007/978-3-0348-0551-3.

\end{thebibliography}

\end{document}
